% ECONOMICS_PROBLEM: {"id": "J12-220", "title": "Partner sorting and inequality", "jel_code": "J12", "source_id": "OP-0367"}
% !TeX program = xelatex
\documentclass[11pt,a4paper]{article}
\usepackage[margin=23mm]{geometry}
\usepackage{fontspec}
\setmainfont{Latin Modern Roman}
\usepackage{amsmath,amssymb,mathtools}
\usepackage{xcolor,enumitem,booktabs,longtable,array,xurl,hyperref}
\definecolor{muted}{HTML}{53636C}
\hypersetup{colorlinks=true,urlcolor=blue,linkcolor=blue}
\setlength{\parindent}{0pt}
\setlength{\parskip}{5pt}
\newcommand{\R}{\mathbb R}
\newcommand{\E}{\mathbb E}
\newcommand{\N}{\mathbb N}
\newcommand{\ind}{\mathbf 1}
\newcommand{\Var}{\operatorname{Var}}
\newcommand{\Cov}{\operatorname{Cov}}
\newcommand{\supp}{\operatorname{supp}}
\DeclareMathOperator*{\argmin}{arg\,min}
\DeclareMathOperator*{\argmax}{arg\,max}
\newcommand{\entrymeta}[3]{{\sffamily\small\color{muted}\textbf{\texttt{#1}}\quad|\quad #2\par #3\par}}
\newcommand{\Problemref}[1]{\texttt{#1}}
\newcommand{\JELref}[2]{\texttt{#2} (JEL \texttt{#1})}
\begin{document}
\section*{Partner sorting and inequality}
\textbf{Problem ID:} \texttt{J12-220}\quad \textbf{JEL:} \texttt{J12}\par
% BEGIN ECONOMICS STATEMENT
\label{problem:OP-0367}
{\sffamily\small\color{muted}Original subject group: Inequality and economic mobility\par}
\label{definitions:problem:30007649}
\textbf{Audit classification:} Background; proposed extension. Published 2026 metadata verified, replacing the older working-paper-only reference. The paper answers exposure and sorting questions; downstream inequality and fertility outcomes are editorial.
\subsection*{Definitions and precise research target}
\textbf{Complete editorial problem.} Fix a cohort of potential partners and a feasible intervention \(z\) changing opportunities to meet across socioeconomic groups. Let \(M(z)\) be the equilibrium matching of partners, allowing some individuals to remain single. Person \(i\) has potential labor earnings \(y_i(M,z)\), reflecting endogenous labor supply as well as partner formation. For household \(h\), define equivalent disposable income \(Y_h(z)\) using a declared sharing and household-size rule. Let \(I(F_z)\) be a specified inequality functional of the household-income distribution \(F_z\), and define the adult-income distribution among children born in a declared calendar window under each policy. Report fertility and cohort-composition changes; different counterfactual children are not identified as the same individual. Estimate \(I(F_z)-I(F_0)\) and child mobility effects. Separate meeting opportunities from partner preferences by varying exposure without assuming preferences remain unchanged after exposure. Distinguish a mechanical reassignment exercise holding earnings fixed from a causal counterfactual that changes work, fertility, and investment. Identification must address selection into institutions, migration, marriage, and parenthood. The proposed problem is to quantify these downstream distributional consequences; evidence that segregation affects who marries whom does not automatically identify the full inequality and next-generation response.

\textbf{Definitions and admissible scope.} A matching \(M\) is a partition of a baseline adult population into admissible pairs and singletons, with no person assigned twice. Household disposable income includes endogenous labor supply, taxes, and a declared equivalence scale. The distribution \(F_z\) is weighted either by persons or by households, with the choice fixed; a Gini requires nonnegative income and positive finite mean. If fertility or parenthood changes, do not label different counterfactual children as the same individual. Define next-generation outcomes as population distributions among births in a declared calendar window under each policy, accompanied by birth counts and parental-composition changes. A fixed-child causal contrast instead requires a baseline cohort and a policy whose relevant population membership is fixed. Choose a common adult-income reference distribution and a child follow-up age. A next-generation mobility target can be the distribution of fixed-reference adult ranks among births in the declared window, conditional on their parents' prereform income group; denote it \(F^{\rm child}_{z,g}\), and compare it with \(F^{\rm child}_{0,g}\) together with group-specific birth counts. This is a policy-specific birth-population comparison, not a same-child causal effect. For the household-inequality target, specify household or person weights and choose \(I\); a nonnegative-income Gini requires \(E[Y_h]>0\) and finite absolute income. Matching feasibility, meeting exposure, labor-supply responses, tax rules, and fertility are fixed primitives of the admissible policy model.
\subsection*{Variants and assumption changes}
\begin{enumerate}
\item Mechanical re-pairing benchmark (editorial): fix each adult's earnings and fertility and change only the matching partition. Compute distributional accounting effects, explicitly without claiming a causal behavioral response.
\item Exposure policy with behavior (editorial extension): change feasible meetings and recompute matching, work, fertility, and child investments in a declared model. The cited marriage paper already estimates exposure effects; downstream inequality and child-cohort consequences require additional identification.
\end{enumerate}
\subsection*{References and source audit}
\textbf{Support and access.} Published 2026 metadata verified, replacing the older working-paper-only reference. The paper answers exposure and sorting questions; downstream inequality and fertility outcomes are editorial. \textbf{Locator:} Published abstract and citation, AER 116(10), pp. 3623--3673.
\begin{enumerate}\raggedright
\item\hypertarget{definitions:source:30007649:1}{} Benjamin Goldman; Jamie Gracie; Sonya R. Porter. 2026. \emph{Who Marries Whom? The Role of Segregation by Race and Class}. Published abstract and citation; American Economic Review 116(10), 3623--3673.
\url{https://swlb1.aeaweb.org/articles?id=10.1257/aer.20250115}.
DOI: \url{https://doi.org/10.1257/aer.20250115}.
\end{enumerate}

\subsection*{Common conventions: Source mathematical conventions}

These conventions apply to each entry unless explicitly replaced.

\textbf{Models and identification.} A primitive model \(m\) specifies all probability laws, feasible choices, information, equilibrium and measurement rules used by its target. The admissible class \(\mathcal M\) is fixed independently of outcomes. If \(P_O(m)\) is its observable probability law and \(\tau(m)\) the target, its identified set at observable law \(P\) is
\[
 \mathcal I_\tau(P)=\{\tau(m):m\in\mathcal M,\ P_O(m)=P\}.
\]
Point identification means that this set is a singleton. An empty set signals incompatibility of data and assumptions. Coordinate bounds need not describe the joint set. Bounds are sharp only if every retained target is attainable or arbitrarily approachable by compatible models. Finite bounds require explicit restrictions on unobserved quantities; unrestricted confounding, hidden wealth, extrapolation or arbitrary equilibrium selection can yield uninformative sets.

\textbf{Causal interventions.} A potential outcome \(Y_i(p)\) is the outcome under a complete policy regime \(p\), including timing, financing, information and market spillovers. The target population and units are fixed before treatment. With interference, use the complete assignment vector or a justified exposure mapping; individual no-interference notation alone cannot identify an economy-wide intervention. A difference of mean potential outcomes does not require the cross-world joint distribution. Conditioning on post-treatment migration, survival, enrollment or employment changes the estimand and needs separate justification.

\textbf{Statistical statements.} An estimator, confidence set, forecast or test specifies a sampling law, model class and loss/coverage criterion. Uniform \((1-\alpha)\)-coverage means \(\inf_{m\in\mathcal M}\Pr_m\{\tau(m)\in C_n\}\geq1-\alpha\), or an explicitly asymptotic version. The whole parameter space trivially has coverage, so informativeness requires a stated width, risk or power objective. A forecasting improvement is relative to a named benchmark, information set and evaluation population.

\textbf{Equilibrium and optimization.} An equilibrium comprises feasible plans, optimality, beliefs where needed, budget constraints and all market/resource clearing. The primitives or a stated benchmark define each correspondence \(\mathcal E(p;m)\). Nonemptiness is assumed or proved, never inferred just from compact policy sets. A selected-equilibrium target uses a declared selection rule; a robust target quantifies over all permitted equilibria. Use suprema or infima unless attainment is established. An \(\varepsilon\)-optimal policy is feasible and within \(\varepsilon\) of the finite optimum value. Report an empty feasible set separately.

\textbf{Welfare and accounting.} Normative utility, interpersonal normalization, social weights, numeraire, discounting and the use of public receipts are primitives. Behavioral utility need not coincide with the welfare criterion. Household welfare includes ownership income and rents once; adding profits again to compensating variation generally double-counts them. A monetary equivalent is defined by an explicit transfer and price convention, with monotonicity and range conditions. Average effects, mechanism contributions, transfers and welfare losses are different targets. Infinite sums require convergence and dynamic feasible plans require terminal or no-Ponzi conditions.

\textbf{Source versions.} A working-paper date, revision date and journal publication year are distinguished. A DOI for a journal version is not automatically the DOI of an earlier manuscript. Locators refer to the stated version and printed pages where specified. A metadata-only check does not verify a claim about a full-text conclusion. Every entry's source subsection narrows the provenance claim accordingly.
% END ECONOMICS STATEMENT
\end{document}
